\section*{Chapitre \ref{ch:m2:reading-a-macro-calendar-and-a-rates-screen} --- Lire un calendrier macroéconomique et un écran de taux}

\begin{solution}{exo:m2:reading-a-macro-calendar-and-a-rates-screen:1}
De minuit le samedi 17 octobre jusqu'à 23 h 59 le jeudi 29 octobre 2026.
\end{solution}

\begin{solution}{exo:m2:reading-a-macro-calendar-and-a-rates-screen:2}
\omterm{def:m2:reading-a-macro-calendar-and-a-rates-screen:bull}{Bull steepening}~: les rendements ont baissé, le 2 ans davantage, si bien que le 2s10s a progressé de 6 points de base.
\end{solution}

\begin{solution}{exo:m2:reading-a-macro-calendar-and-a-rates-screen:3}
2s10s $415 - 361 = 54$, 5s30s $477 - 371 = 106$, \omterm{def:m2:reading-a-macro-calendar-and-a-rates-screen:steep}{butterfly de courbe} $2 \times 371 - 361 - 415 = -34$ points
de base.
\end{solution}

\begin{solution}{exo:m2:reading-a-macro-calendar-and-a-rates-screen:4}
Le rendement à 2 ans est proche du taux directeur moyen attendu sur deux ans, et le
rapport modifie cette anticipation de la manière la plus directe~; le rendement à 10 ans intègre aussi des taux attendus
plus lointains, qui bougent moins avec les données d'un seul mois, et une prime de terme.
\end{solution}

\begin{solution}{exo:m2:reading-a-macro-calendar-and-a-rates-screen:5}
$50\,000 \times (-5 + 15) = \text{USD}~500\,000$~; une baisse parallèle de 10 ne rapporte rien au
premier ordre.
\end{solution}

\begin{solution}{exo:m2:reading-a-macro-calendar-and-a-rates-screen:6}
Comme le résultat est proche d'un pile ou face pour quiconque n'a pas d'edge sur le chiffre, les
mouvements sont deux à trois fois plus amples que d'habitude, et la liquidité se raréfie autour de la publication~:
le risque croît davantage que la récompense espérée.
\end{solution}

\begin{solution}{exo:m2:reading-a-macro-calendar-and-a-rates-screen:7}
L'erreur type est d'environ le bruit divisé par la racine carrée du nombre de publications,
$3/\sqrt n$ avec des \omterm{def:m2:reading-a-macro-calendar-and-a-rates-screen:surprise}{surprises statistiques} de variance unitaire~: 0.5 requiert environ 36 publications, soit trois ans de
données mensuelles. La simulation avec 36 publications donne une erreur type moyenne de 0.51.
\end{solution}

\begin{solution}{exo:m2:reading-a-macro-calendar-and-a-rates-screen:8}
Les 10 points de base représentent le mouvement absolu moyen du rendement à 2 ans, et non une mauvaise journée, et ils
ne disent rien des autres expositions de la position~: le risque d'une position de courbe est le mouvement de la pente
multiplié par son DV01, celui d'un portefeuille étant la combinaison de ses expositions le long de la courbe. Utiliser
la distribution des mouvements des jours d'événement pour chaque exposition, y compris sa queue~: le
rendement à 2 ans a varié jusqu'à 30 points de base les jours de payrolls dans l'échantillon.
\end{solution}

\begin{solution}{pb:m2:reading-a-macro-calendar-and-a-rates-screen:1}
\textbf{1.} 0.0183 pour le 2 ans et 0.0772 pour le 10 ans, par 100.
\textbf{2.} 77\,168~USD par point de base.
\textbf{3.} 421~millions de dollars de notionnel sur le 2 ans.
\textbf{4.} Parce que le 10 ans fait varier le prix quatre fois plus par point de base~:
à notionnel égal, il s'agirait surtout d'un pari sur le niveau du 10 ans.
\textbf{5.} Rien, au premier ordre.
\textbf{6.} 694\,512~USD au premier ordre~; 688\,510~USD après réévaluation complète.
\textbf{7.} Une perte de 771\,680~USD au premier ordre~; 772\,969~USD après réévaluation.
\textbf{8.} \omterm{def:m2:reading-a-macro-calendar-and-a-rates-screen:bull}{Bull steepening}, puis \omterm{def:m2:reading-a-macro-calendar-and-a-rates-screen:bull}{bear flattening}.
\textbf{9.} 34 fois sur 44.
\textbf{10.} La convexité et les différences d'ampleur des mouvements~: les prix du 2 ans et du 10 ans
ne sont pas linéaires par rapport à leurs rendements, et la convexité du 10 ans est plus grande.
\textbf{11.} Non~: elle commence le lendemain, samedi 5 septembre.
\textbf{12.} Les décideurs monétaires ne peuvent pas le commenter avant la réunion, si bien que la lecture
du CPI par le marché n'est ni corrigée ni confirmée par eux avant la décision.
\textbf{13.} 10.2 points de base contre 4.6, soit 2.2 fois autant.
\textbf{14.} Le CPI du 11 septembre, les adjudications du Trésor et tout autre événement programmé, comme
par exemple les discours prononcés avant le début de la \omterm{def:m2:reading-a-macro-calendar-and-a-rates-screen:blackout}{période de silence}.
\textbf{15.} Le plus grand mouvement d'un jour de payrolls sur la pente 2s10s dans l'échantillon a été de 11 points
de base (1\up{er} août 2025)~: un DV01 d'environ 45\,000~USD par jambe maintient cette perte dans
la limite de 500\,000~USD~; la position du trader, de 77\,000~USD par point de base, perdrait environ 850\,000~USD un tel jour, face à ce plafond.
\textbf{16.} Parce qu'un rapport faible peut faire bouger la partie longue de la courbe tout autant, s'il soulève des doutes sur
la prime de terme ou l'offre, ou parce que le marché l'avait déjà intégré~; et parce que le
rapport n'est qu'un des éléments qui ont fait bouger les rendements ce jour-là.
\textbf{17.} Le 2 ans, la \omterm{def:m2:short-term-interest-rate-futures:path}{trajectoire implicite des taux directeurs} et le 2s10s avant et après~; la réaction
de la partie longue~; les volumes et la liquidité~; les prochains événements du calendrier.
\textbf{18.} Le \omterm{def:m2:reading-a-macro-calendar-and-a-rates-screen:steep}{steepener} rapporte lorsque le marché abaisse sa trajectoire attendue des taux directeurs
à court terme~: c'est une façon de parier sur la trajectoire du chapitre 8 sans détenir le niveau des taux.
\textbf{19.} Résultat nommé~: \emph{vendredi des payrolls}~: acheter 421~millions de dollars de notionnel sur le 2 ans
contre 100~millions de dollars sur le 10 ans, soit 77\,168~USD par point de base sur chaque jambe~; le
\omterm{def:m2:reading-a-macro-calendar-and-a-rates-screen:steep}{steepener} gagne 694\,512~USD sur une répétition du 2 août 2024 (688\,510~USD après réévaluation) et perd
771\,680~USD sur une répétition du 4 octobre 2024 (772\,969~USD après réévaluation).
\textbf{20.} Sur la pente entre les deux échéances, et non sur le niveau des taux.
\end{solution}

\begin{solution}{iq:m2:reading-a-macro-calendar-and-a-rates-screen:1}
Le rapport sur l'emploi et le CPI, à 8 h 30, heure de l'Est, le premier ayant lieu
généralement le premier vendredi du mois et le second vers le milieu~; les décisions du FOMC après ses huit
réunions prévues~; s'ajoutent les ventes au détail, le PIB, l'inflation PCE, les refinancements du Trésor et
les adjudications.
\iqlookfor{les principales publications et leur calendrier.}
\end{solution}

\begin{solution}{iq:m2:reading-a-macro-calendar-and-a-rates-screen:2}
\omterm{def:m2:reading-a-macro-calendar-and-a-rates-screen:bull}{Bull steepening}~: les rendements baissent, le court terme davantage, comme lorsque des données faibles tirent
à la baisse les taux directeurs attendus. \omterm{def:m2:reading-a-macro-calendar-and-a-rates-screen:bull}{Bear flattening}~: les rendements montent, le court terme davantage, comme lorsque des données fortes ou une
banque centrale aux positions restrictives pousse les taux directeurs attendus à la hausse.
\iqlookfor{sens, forme et canal de la trajectoire des taux directeurs.}
\end{solution}

\begin{solution}{iq:m2:reading-a-macro-calendar-and-a-rates-screen:3}
Acheter le 2 ans et vendre le 10 ans à DV01 égaux, en obligations, en contrats à terme ou en swaps. Elle
gagne lorsque le 2s10s s'élargit. Risques~: la pente évoluant dans l'autre sens, les mouvements non parallèles au sein
de chaque jambe, le financement et le roll (le carry de chaque jambe), et la base entre les instruments.
\iqlookfor{neutralité au DV01 et risques restants.}
\end{solution}

\begin{solution}{iq:m2:reading-a-macro-calendar-and-a-rates-screen:4}
Recueillir la publication, le consensus et le mouvement du marché dans une fenêtre étroite autour de
l'heure de publication~; standardiser la \omterm{def:m2:reading-a-macro-calendar-and-a-rates-screen:surprise}{surprise statistique}~; régresser les mouvements sur les surprises sur un
grand nombre de publications~; vérifier la stabilité du coefficient dans le temps et si les grandes surprises
agissent de façon linéaire.
\iqlookfor{\omterm{def:m2:reading-a-macro-calendar-and-a-rates-screen:surprise}{surprise statistique}, fenêtre étroite, régression, stabilité.}
\end{solution}

\begin{solution}{iq:m2:reading-a-macro-calendar-and-a-rates-screen:5}
Les niveaux et variations de la courbe aux échéances clés, les pentes et un \omterm{def:m2:reading-a-macro-calendar-and-a-rates-screen:steep}{butterfly de courbe}~; la \omterm{def:m2:short-term-interest-rate-futures:path}{trajectoire implicite des taux directeurs}
et sa variation~; les \omterm{def:m2:interest-rate-swaps:spread}{spreads de swap} et la \omterm{def:m2:fx-swaps-forwards-and-the-cross-currency-basis:basis}{base de devises}~; la volatilité implicite~; les taux de financement~;
le calendrier de la journée et le DV01 des positions par tranche face à leurs limites.
\iqlookfor{niveau, forme, anticipations, valeur relative, risque, calendrier.}
\end{solution}

\begin{solution}{iq:m2:reading-a-macro-calendar-and-a-rates-screen:6}
S'abonner au flux des publications avec la latence la plus faible disponible et analyser les chiffres
dans des champs standard~; calculer les surprises par rapport au consensus préchargé~; transmettre les nouveaux
chiffres aux moteurs de valorisation, qui recalculent de façon incrémentale les courbes et les cotations~; les contrôles de
limites et la mise à jour du risque dans la même boucle d'événements~; tout horodaté pour le rejeu~; des
coupe-circuits si le flux est mal formé ou en retard.
\iqlookfor{contexte préchargé, recalcul incrémental, sécurité intégrée.}
\end{solution}
